\newpage
\spread{Learning what changes a translation}
\noindent Errors in two visual features can have the same magnitude and different consequences for a sentence. We continued the trained optical reader under two objectives and measured what each reader let a frozen bilingual decoder say.
\platehere{Figure_216_prediction_sensitive_reading}{Training against the decisions of a frozen decoder}{\textbf{a}, Full-vocabulary KL divergence from the original translation model on two held development books, at 3,191 English and 3,197 Chinese target decisions. Lower is better. \textbf{b}, Complete meanings from 96 fresh passages translated in both directions, across four previously assessed books. The automated judge received the source and candidate answers in both orders. The first four methods use a pretrained piece decoder. The independent ILM emits glyph images. Qwen3-8B Q4 receives text.}
\begin{multicols}{2}
\section*{Teaching one decision at a time}
Both readers received four more passes over the same 28,342 pictured training pairs in identical order. Both regressed the contextual source vectors and attention weights. One also matched the teacher's next-piece distribution at every genuine target prefix, including the stop. We retained the 32 highest-probability teacher pieces and combined the rest into an event $R$. If $P$ and $Q$ denote full teacher and student distributions and $G$ groups these events, the chain rule gives
\[
D(P\Vert Q)=D(GP\Vert GQ)+P(R)D(P_{|R}\Vert Q_{|R}).
\]
The grouped term thus measures part of the full divergence. Every output still contributes to the student's normalizer. The 32 pieces held 84.9\% of teacher mass on English development decisions and 79.4\% on Chinese decisions.
\section*{What changed in free writing}
For the Chinese sentence meaning \emph{He hesitated}, the regression reader yielded \emph{He was wondering.} With prediction loss the same pretrained decoder produced \emph{He hesitated.} On the first request, which described a broad hand and concealed dislike, the latter readout lost the gesture and the concealment. Both examples are text readouts of an optical source field.
\par Prediction training earned 22/192 complete meanings against 10 for continued regression and 11 at the starting checkpoint. Most of the gain was into English (15 versus four). Its advantage across the four works was $+3,+8,-1,+2$ answers. The paired work-bootstrap 95\% interval was 0 to 13.54 percentage points, missing the prespecified positive lower bound. In a separate shorter candidate list the same predictor earned 33 credits versus two after source images were exchanged. The different counts for identical answers show the sensitivity of automatic grading to the comparison set.
\end{multicols}
\newpage
\spread{Can a compact translator write its own pixels?}
\noindent The 1.69-million-parameter writer's choice of the next glyph carried linguistic information, but its stored drawings supplied every stroke. We retained its visual choices and taught a shared neural network to turn the selected 160-dimensional optical vector into 1,024 binary pixel decisions.
\platehere{Figure_217_neural_pixel_output}{Full answers from a learned glyph motor}{Both outputs were generated on CPU by the released image-to-image model. The only input to the run was a pictured sentence, its occupied cell count and the requested output language. \textbf{a}, First held answer retaining a whole-meaning credit. The writer translated the speaker's agreement. \textbf{b}, First rejected answer in assessment order. The repeated words and damaged meaning come from the model. These original PNGs retain their full cell spacing and aspect ratio.}
\begin{multicols}{2}
\section*{Writing and reading its own strokes}
Let $c_t$ be the visual choice made by the language network at step $t$. The motor predicts $I_t=\mathbf{1}[M_\phi(c_t)>0]$. The writer pools $I_t$ over $2\times2$ blocks to form the optical input to its next step. The deployed bundle contains trained weights and optical choice vectors. It uses no glyph bitmap table, font renderer or pretrained text decoder. The set of codes is finite, which limits the forms it can express.
\par We fitted the motor to 3,717 distinct drawings in the writer's repertoire. At the final weights it recovered all 3,806,208 training pixels. This measures memory for the known repertoire. The new motor had 2,266,112 trainable parameters. The deployed model has 3,954,592 parameters and 604,480 stored optical-code scalars. The four thousand motor updates took 4.2 seconds on one RTX 4090~D after loading the antecedent models. This is the cost of this final fit only.
\section*{A check on the whole sequence}
If two motors produce the same bitmap on each choice visited by a deterministic writer starting at the same state, their next inputs, complete outputs and stop decisions agree by induction. A Lean proof includes both natural stop and the finite generation limit. Floating-point implementations still require a test. On 192 fresh book requests, the network's output PNG, cell count and stop matched the repertoire writer on all 192. Two replies earned complete-meaning credits on the source-only assessment. For comparison the text baseline earned 160. Learning the strokes preserved the strengths and the linguistic errors of the original writer.
\end{multicols}
\noindent{\small \href{https://github.com/lachlanchen/ImagizedLanguageModel/blob/main/publication/ilm-nature/evidence/neural_motor_answers.zip}{All 192 pictured inputs and raw output PNGs}. \href{https://github.com/lachlanchen/ImagizedLanguageModel/blob/main/publication/ilm-nature/evidence/visual_output_figure_index.json}{Figure-by-figure evidence index}.}
